% ===== PRÉAMBULE CANONIQUE — à coller VERBATIM en tête de CHAQUE .tex =====
\documentclass[11pt,a4paper]{article}
\usepackage[utf8]{inputenc}\usepackage[T1]{fontenc}\usepackage{lmodern}
\usepackage[french]{babel}
\usepackage{amsmath,amssymb,mathtools}\usepackage{icomma}
\usepackage{siunitx}\sisetup{locale=FR}  % \qty{}{}, \si{}, \ang{}
\usepackage{xcolor}\usepackage{colortbl}
\usepackage{tikz}\usepackage{pgfplots}\pgfplotsset{compat=1.18}
\usepackage[siunitx,europeanresistors]{circuitikz} % circuits (élec.)
\usepackage[most]{tcolorbox}\usepackage{enumitem}\usepackage{array,booktabs}
\usepackage[a4paper,margin=2.2cm]{geometry}
\usetikzlibrary{arrows.meta,calc,angles,quotes,positioning,patterns,%
  backgrounds,shapes.geometric,decorations.pathmorphing,decorations.markings,intersections}
\definecolor{primary}{RGB}{222,49,99}\definecolor{primarydark}{RGB}{150,22,55}
\definecolor{defcol}{RGB}{0,114,178}\definecolor{methcol}{RGB}{52,92,148}
\definecolor{excol}{RGB}{0,135,90}\definecolor{attncol}{RGB}{200,40,55}
\tcbset{boxbase/.style={enhanced,breakable,boxrule=0.9pt,arc=2.5pt,
  left=3mm,right=3mm,top=2.3mm,bottom=2.3mm,fonttitle=\bfseries,coltitle=white}}
% --- boîtes TUTORIEL (noms EXACTS, ne pas renommer) ---
\newtcolorbox{cle}[1][]{boxbase,colback=defcol!6,colframe=defcol,
  title={\textsc{À retenir}\ifx\relax#1\relax\relax\else\textmd{ : \itshape #1}\fi}}
\newtcolorbox{methode}[1][]{boxbase,colback=methcol!7,colframe=methcol,sharp corners,
  title={\textsc{Méthode}\ifx\relax#1\relax\relax\else\textmd{ --- \itshape #1}\fi}}
\newtcolorbox{exemplebox}[1][]{boxbase,colback=excol!5,colframe=excol,
  title={\textsc{Exemple}\ifx\relax#1\relax\relax\else\textmd{ : \itshape #1}\fi}}
\newtcolorbox{piege}[1][]{boxbase,colback=attncol!6,colframe=attncol,
  title={\textsc{Piège}\ifx\relax#1\relax\relax\else\textmd{ : \itshape #1}\fi}}
\newtcolorbox{remarque}[1][]{boxbase,colback=black!4,colframe=black!45,
  title={\textsc{Remarque}\ifx\relax#1\relax\relax\else\textmd{ : \itshape #1}\fi}}
% --- boîtes EXERCICE / CORRIGÉ (noms EXACTS, auto-comptées) ---
\newtcolorbox[auto counter]{exo}[1][]{boxbase,colback=primary!4,colframe=primary,
  title={\textsc{Exercice}~\thetcbcounter\ifx\relax#1\relax\relax\else\textmd{ --- \itshape #1}\fi}}
\newtcolorbox[auto counter]{corrige}[1][]{boxbase,colback=excol!4,colframe=excol,
  title={\textsc{Corrigé}~\thetcbcounter\ifx\relax#1\relax\relax\else\textmd{ --- \itshape #1}\fi}}
\newcommand{\vv}[1]{\vec{#1}}\newcommand{\ds}{\displaystyle}
\newcommand{\R}{\mathbb{R}}\newcommand{\N}{\mathbb{N}}\newcommand{\Z}{\mathbb{Z}}
% (un \usepackage{fancyhdr}+en-tête est possible ; il est ignoré côté web)
% =========================================================================
\begin{document}
\sisetup{per-mode=symbol}
Sauf mention contraire, $g=\qty{10}{\metre\per\second\squared}$ et frottements négligés.

\begin{exo}[hauteur maximale avec vitesse initiale]
\begin{enumerate}[label=\alph*)]
  \item Un sac de sable est lâché d'une montgolfière située à \qty{120}{\metre} du sol qui \emph{monte}
        à la vitesse constante de \qty{10}{\metre\per\second}. Quelle hauteur maximale (par rapport au sol) le sac atteint-il ?
  \item Une balle en chute libre rebondit sur une plateforme située à \qty{3}{\metre} de hauteur ; après
        l'impact, sa vitesse ascensionnelle vaut \qty{10}{\metre\per\second}. Quelle hauteur maximale atteint-elle par rapport au sol ?
\end{enumerate}
\end{exo}

\begin{exo}[quand la trajectoire fait 45°]
Une bille est projetée horizontalement du sommet d'une tour à $v_{0x}=\qty{40}{\metre\per\second}$.
\begin{enumerate}[label=\alph*)]
  \item Après combien de temps le vecteur vitesse fait-il un angle de \ang{45} avec l'horizontale ?
        (Indice : cela se produit quand $v_y=v_x$.)
  \item De quelle hauteur la bille est-elle alors descendue, et quelle est la norme de sa vitesse à cet instant ?
\end{enumerate}
\end{exo}

\begin{exo}[la masse n'y change rien]
\begin{enumerate}[label=\alph*)]
  \item Deux masses $m$ et $2m$ sont lancées du sol, en même temps et à la même vitesse verticale.
        Atteignent-elles la même hauteur ? Au même instant ? Justifie sans calcul.
  \item Sur la Lune ($g=\qty{1,6}{\metre\per\second\squared}$), on lance une masse verticalement à \qty{20}{\metre\per\second}. Quelle hauteur maximale atteint-elle ?
\end{enumerate}
\end{exo}

\begin{exo}[la profondeur du puits]
Un randonneur lâche une pierre dans un puits. Le bruit de l'impact lui parvient \qty{4,25}{\second}
après le lâcher. La vitesse du son est \qty{320}{\metre\per\second}.
\begin{enumerate}[label=\alph*)]
  \item Exprime la durée totale comme la somme du temps de chute $t_c$ et du temps de remontée du son $t_s$.
  \item Vérifie que $t_c=\qty{4}{\second}$ convient et déduis-en la profondeur du puits.
\end{enumerate}
\end{exo}

\begin{exo}[lancer vertical : vrai ou faux]
Une balle est lancée du sol verticalement vers le haut à \qty{10}{\metre\per\second}. Indique si chaque
affirmation est vraie ou fausse, puis compte les vraies.
\begin{enumerate}[label=\alph*)]
  \item Elle atteint une hauteur maximale de \qty{5}{\metre}.
  \item Elle retombe au sol \qty{1}{\second} après le lancer.
  \item Sa vitesse est nulle au sommet de la trajectoire.
  \item Elle repasse au sol avec une vitesse de \qty{10}{\metre\per\second}.
\end{enumerate}
\end{exo}

\end{document}
